\section{Požadavky na systém}
\label{sec:requirements}

Před návrhem systému je nutné vymezit, \emph{co} má systém poskytovat (funkční požadavky) a~\emph{jaké vlastnosti} má při provozu splňovat (nefunkční požadavky).

\subsection{Funkční požadavky}
\label{subsec:requirements:functional}

\begin{enumerate}
  \item \textbf{Webová aplikace s~jednoduchým přístupem}
    \begin{itemize}
      \item Aplikace je dostupná v~prohlížeči, bez nutnosti instalace klientského software.
      \item Aplikace je read-only a~nevyžaduje autentizaci uživatelů, aby ji mohl používat kdokoli ze zájemců o~vývoj sledovaného repozitáře.
    \end{itemize}

  \item \textbf{Vždy aktuální data o~sledovaném repozitáři}
    \begin{itemize}
      \item Aplikace průběžně získává nová data o~\PR, issues a~aktivitě nad soubory.
      \item Aplikace umožňuje opakované spouštění synchronizace bez ztráty konzistence dat a~doplňuje nové i~změněné záznamy.
    \end{itemize}

  \item \textbf{Historizace stavů v~čase}
    \begin{itemize}
      \item Aplikace uchovává časový vývoj stavů a~atributů \PR a~issues, aby bylo možné rekonstruovat jejich historický stav.
      \item Aplikace dokáže historii doplnit (\emph{backfill}) i~zpětně pro záznamy vzniklé v~době, kdy nebyla v~provozu.
    \end{itemize}

  \item \textbf{Analytické pohledy a~filtrování}
    \begin{itemize}
      \item Aplikace umožňuje zobrazovat statistiky a~dotazy v~různých časových oknech.
      \item Aplikace umožňuje filtrovat výsledky podle relevantních parametrů (např. týmu, cesty v~repozitáři, identifikátoru \PR či issue).
      \item Pro práci s~většími datovými sadami je k~dispozici stránkování výsledků.
      \item Konkrétní sada poskytovaných dotazů je vymezena \customref{v~Podsekci}{subsec:requirements:queries}.
    \end{itemize}

  \item \textbf{Self-healing po výpadku}
    \begin{itemize}
      \item Pokud aplikace nějakou dobu neběží nebo dojde k~chybě při synchronizaci, je schopna chybějící data doplnit při dalším spuštění.
      \item Aplikace minimalizuje duplicitní záznamy a~udržuje konzistentní stav i~při opakovaném zpracování stejných vstupů.
    \end{itemize}

  \item \textbf{Jednoduchý deployment}
    \begin{itemize}
      \item Aplikaci je možné spustit v~lokálním i~serverovém prostředí bez náročné ruční konfigurace.
      \item Citlivé údaje, jako například přístupové tokeny, jsou předávány externí konfigurací, nikoliv pevně zapsány v~kódu.
    \end{itemize}
\end{enumerate}

\subsection{Specifikace analytických dotazů}
\label{subsec:requirements:queries}

Funkční požadavek na analytické pohledy vyžaduje sadu konkrétních dotazů nad nasbíranými daty. Tato podsekce pro každý dotaz uvádí jeho účel, vstupní parametry a~očekávaný výstup. Návrh realizace jednotlivých dotazů je popsán v~\customref{Sekci}{sec:queries}.

\begin{description}
  \item[Q1 -- Historický stav \PR k~danému datu]
    Při revizi kódu ve velkém repozitáři je užitečné zpětně zjistit, v~jakém stavu se konkrétní \PR nacházel v~určitém období. Tento dotaz umožňuje rekonstruovat historii stavů~\PR na základě zaznamenaných událostí a~změn štítků.

    \textbf{Vstup:} identifikátor~\PR, počáteční datum, počet dní zpětně.\\
    \textbf{Výstup:} seznam stavů~\PR v~zadaném časovém rozmezí, každý s~časovým razítkem.

  \item[Q2 -- Počet \PR v~konkrétním stavu k~datu]
    Pro sledování trendů v~procesu revize kódu je potřeba zjistit, kolik~\PR se nacházelo v~daném stavu k~určitému dni. Opakovaným voláním tohoto dotazu pro různá data lze sestavit časovou řadu zachycující vývoj počtu~\PR v~konkrétním stavu.

    \textbf{Vstup:} stav~\PR (výčtový typ), datum.\\
    \textbf{Výstup:} počet~\PR v~zadaném stavu k~danému datu.

  \item[Q3 -- \PR čekající na zpracování]
    \label{item:query:q3}
    V~rozsáhlém repozitáři snadno vzniká situace, kdy některé~\PR dlouho čekají na pozornost. Tento dotaz zobrazuje aktuální~\PR ve stavech \code{S-waiting-on-review}, \code{S-waiting-on-bors} nebo \code{S-waiting-on-author} \sectionref{sec:rust_labels}, čímž umožňuje identifikovat nejstarší čekající~\PR.

    \textbf{Vstup:} (žádné specifické parametry, pracuje s~aktuálním stavem).\\
    \textbf{Výstup:} stránkovaný seznam~\PR seřazený podle doby čekání.

  \item[Q4 -- Události konkrétní issue k~datu]
    Pro pochopení vývoje konkrétní issue je potřeba vidět chronologický přehled jejích událostí. Dotaz vrací zaznamenané akce jako uzavření, znovuotevření nebo komentáře k~zadanému datu.

    \textbf{Vstup:} identifikátor issue, datum.\\
    \textbf{Výstup:} seznam událostí dané issue k~zadanému datu s~časovými razítky. Pokud issue neexistuje, vrací se prázdná odpověď.

  \item[Q5 -- Nejčastěji upravované soubory přispěvatelem]
    Při rozdělování práce v~týmu je užitečné vědět, na kterých částech projektu se daný vývojář nejvíce podílí. Dotaz vrací soubory, které přispěvatel nejčastěji upravoval v~zadaném období.

    \textbf{Vstup:} jméno přispěvatele, počet dní zpětně od aktuálního data.\\
    \textbf{Výstup:} seřazený seznam souborů s~počtem úprav a~metadaty přispěvatele.

  \item[Q6 -- Přispěvatelé, kteří upravili soubor nebo složku]
    Tento dotaz nabízí opačný pohled k~Q5. K~zadané cestě souboru nebo adresáře vrací přispěvatele, kteří v~dané oblasti prováděli změny. Cesta je vyhledávána jako prefix, takže dotaz pokrývá i~celý adresář včetně podadresářů.

    \textbf{Vstup:} cesta k~souboru nebo složce, časové okno.\\
    \textbf{Výstup:} stránkovaný seznam přispěvatelů s~metadaty a~celkovým počtem.

  \item[Q7 -- Soubory upravené členy týmu]
    Pro koordinaci práce mezi týmy je potřeba zjistit, na kterých souborech nebo modulech pracují členové konkrétního týmu. Výsledek lze zobrazit jako plochý seznam, stromovou strukturu seskupenou do zvolené hloubky adresářové hierarchie, nebo jako úplný adresářový strom.

    \textbf{Vstup:} název týmu, časové okno, počet úrovní hierarchie pro seskupení.\\
    \textbf{Výstup:} seznam souborů (plochý nebo hierarchický) seřazený podle počtu úprav.

\end{description}

\subsection{Nefunkční požadavky}
\label{subsec:requirements:nonfunctional}

\begin{enumerate}
  \item \textbf{Výkon}
    \begin{itemize}
      \item Aplikace musí efektivně zpracovávat data z~rozsáhlého repozitáře s~velkým množstvím \PR a~issues.
      \item Odezva analytických dotazů musí zůstat únosná i~nad historickými daty v~řádu desítek tisíc \PR.
    \end{itemize}

  \item \textbf{Spolehlivost a~odolnost}
    \begin{itemize}
      \item Aplikace musí korektně zpracovat i~neúplná data z~externích zdrojů.
      \item Při dílčích chybách musí být chyba zaznamenána a~musí být možné synchronizaci zopakovat bez poškození již uložených dat.
    \end{itemize}

  \item \textbf{Konzistence a~kvalita dat}
    \begin{itemize}
      \item Aplikace musí minimalizovat duplicitní záznamy a~udržovat konzistentní stav.
      \item Aplikace musí uchovávat časové souvislosti událostí pro pozdější analýzu trendů.
    \end{itemize}

  \item \textbf{Provozovatelnost}
    \begin{itemize}
      \item Aplikace musí být konfigurovatelná pomocí externí konfigurace (sledovaný repozitář, přístupové tokeny, připojení k~databázi).
      \item Aplikace musí poskytovat logování pro diagnostiku běhu a~chyb.
    \end{itemize}

  \item \textbf{Bezpečnost}
    \begin{itemize}
      \item Přístupové tokeny k~externím službám nesmí být pevně zapsány v~kódu.
      \item Vstupní parametry uživatelských dotazů musí být validovány a~chybové stavy bezpečně zpracovány.
    \end{itemize}
\end{enumerate}
